\documentclass[11pt]{amsart}
\usepackage[T1]{fontenc}
\usepackage[utf8]{inputenc}
\usepackage{amsmath,amssymb,amsthm,mathtools}
\usepackage{microtype}
\usepackage[colorlinks=true,linkcolor=blue,citecolor=blue,urlcolor=blue]{hyperref}

\newtheorem{theorem}{Theorem}
\newtheorem{lemma}[theorem]{Lemma}

\theoremstyle{remark}
\newtheorem{remark}[theorem]{Remark}

\newcommand{\Cyc}[1]{C_{#1}}
\newcommand{\Z}{\mathbb{Z}}
\newcommand{\ord}{\operatorname{ord}}


\begin{document}

\title[Nonexistence of a Cyclic $(963,222,51)$ Difference Set]{Nonexistence of a Cyclic $(963,222,51)$ Difference Set\thanks{AutoMath assisted in generating the solution and preparing this draft.}}
\date{April 15, 2026}

\begin{abstract}
We show that the cyclic group $\Cyc{963}$ admits no $(963,222,51)$-difference set.
The proof uses the quotient $\Cyc{963}\twoheadrightarrow \Cyc{107}$.
A contracted multiplier theorem makes $19$ a $107$-multiplier, so after translation the quotient coefficient function is constant on the multiplication-by-$19$ orbits on $\Cyc{107}$.
Those orbits have sizes $1,53,53$, while each quotient coefficient is at most $9$.
The required weight $222$ is therefore impossible.
\end{abstract}

\maketitle

\section{Introduction}
The cyclic row $(963,222,51)$, equivalently the $[9,107]$ case, appears as a small open residue in Gordon's La Jolla difference-set status tables \cite{Gordon2019}.
We give a short quotient-orbit obstruction that rules it out.

\begin{theorem}
The cyclic group $\Cyc{963}$ does not admit a $(963,222,51)$-difference set.
\end{theorem}

\begin{proof}
Assume for contradiction that $D\subseteq G=\Cyc{963}$ is a $(963,222,51)$-difference set.
Set
\[
n=k-\lambda=222-51=171=3^2\cdot 19.
\]
Let $N$ be the subgroup of order $9$, and write
\[
Q=G/N\cong \Cyc{107}.
\]
For each $u\in Q$, let $a_u$ denote the number of points of $D$ in the coset above $u$.
Then
\[
0\le a_u\le 9
\qquad\text{and}\qquad
\sum_{u\in Q} a_u = 222.
\]

We use the contracted multiplier theorem on the quotient $Q$.
Take $t=19$ and $w=107$.
For the prime divisor $19$ of $n$ we have $19^1\equiv 19\pmod{107}$, and for the prime divisor $3$ of $n$ we have
\[
3^{95}\equiv 19 \pmod{107}.
\]
Therefore Baumert--Gordon's sufficient criterion for $w$-multipliers applies, so $19$ is a $107$-multiplier for $D$ \cite[Thm.~3.2]{BG2004}.
Equivalently, after translating $D$ if necessary, we may assume that the quotient coefficient function is invariant under multiplication by $19$ on $Q$:
\[
a_{19u}=a_u \qquad (u\in Q).
\]

Since $107$ is prime, the orbit sizes of $u\mapsto 19u$ are determined by $\ord_{107}(19)$.
A direct modular computation gives
\[
\ord_{107}(19)=53.
\]
Hence the action of multiplication by $19$ on $Q$ has exactly three orbits:
\[
\{0\},\qquad O_1,\qquad O_2,
\]
where $|O_1|=|O_2|=53$.
Because the function $u\mapsto a_u$ is orbit-constant, there exist integers $b_0,\alpha,\beta$ with
\[
0\le b_0,\alpha,\beta\le 9
\]
such that
\[
222 = b_0 + 53(\alpha+\beta).
\]
Reducing modulo $53$ gives
\[
b_0 \equiv 222 \equiv 10 \pmod{53}.
\]
But $0\le b_0\le 9$, which is impossible.
This contradiction proves that $\Cyc{963}$ admits no $(963,222,51)$-difference set.
\end{proof}

\begin{remark}
The decisive obstruction takes place on the prime quotient $\Cyc{107}$.
The weight $222$ is too small in the distinguished fixed coset and too rigid modulo the two $53$-cycles.
\end{remark}

\begin{thebibliography}{99}

\bibitem{BG2004}
Leonard D.~Baumert and Daniel M.~Gordon,
\emph{On the existence of cyclic difference sets with small parameters},
Fields Inst. Commun. \textbf{41} (2004), 29--40.

\bibitem{Gordon2019}
Daniel M.~Gordon,
\emph{The La Jolla Difference Set Repository},
ArasuFest talk slides, August 3, 2019.

\end{thebibliography}

\end{document}
